% Limitations section content below
% Note: \section{} header is in main.tex

Our findings are subject to several scope constraints that bound their generalizability.

\paragraph{Model Scale.}
All experiments use Qwen3-0.6B (596M parameters). Larger models produce higher-quality candidates with greater score variance, which could amplify curation effects in either direction: stronger selection pressure from wider score gaps, or greater resilience from richer internal representations. Whether the semantic stability we observe holds at the 7B or 70B scale remains an open question.

\paragraph{Generation Depth.}
We evaluate three generations of iterative self-training. The model collapse literature typically requires more iterations before distributional narrowing becomes severe \cite{placeholder_shumailov}. Our results establish that reward curation is safe over three generations, but they do not guarantee stability at generation ten or beyond. Subtle compounding effects below our measurement threshold could accumulate over longer horizons.

\paragraph{Domain Coverage.}
The UltraFeedback prompt set (2,200 instructions) covers general instruction-following across diverse topics. Narrow domains such as code generation or mathematical reasoning impose tighter output constraints and may exhibit different collapse dynamics. In particular, domains where correct answers cluster tightly in embedding space could reduce the protective effect of prompt diversity.

\paragraph{Single Reward Model.}
We use ArmoRM-Llama3-8B-v0.1 as the sole reward signal throughout all experiments. Different reward models encode different biases: some may penalize structural variety more aggressively, others may exhibit sharper score distributions that strengthen curation pressure. The specific format collapse patterns we document (opening template convergence, near-duplicate starts) reflect ArmoRM's preferences and may not generalize to all reward models.

\paragraph{Format Metric Granularity.}
Opening entropy and near-duplicate start rate capture homogenization in the first 10 tokens and 50 characters of each response, respectively. These metrics are blind to structural convergence in middle sections, concluding paragraphs, or organizational patterns such as list usage and paragraph structure. A complete characterization of format collapse would require finer-grained positional analysis that our current metrics do not provide.